% wrapped: the wrappedfigure environment, a key-value front end for wrapfig's
% wrapfigure. The filter writes
% \begin{wrappedfigure}[side=right,width=0.4\linewidth,gap=1em,lines=12].
%
% The sides are wrapfig's floating placements, R and L: a figure that does
% not fit on the page moves to the start of a paragraph on the next page.
%
% wrapfig puts \columnsep between the figure and the text, and it reads
% \columnsep when it places the figure, at the start of a later paragraph.
% So the environment sets \columnsep to the gap when it ends, and wrapfig's
% \WF@startwrapping and \WFclear put the document's value back once the
% figure is placed or put out without text.
%
% wrapfig holds one figure at a time. So a wrappedfigure first deals with an
% earlier one. A figure that moved and still waits for a paragraph is put
% out here, centered, with no text beside it. A wrap whose text is still
% narrowed ends: the environment skips the lines still left beside the
% earlier figure, and \WFclear ends the wrap. \WFclear ends it inside the
% environment's group, so \wrapped@afterenv ends it again after the group.
\usepackage{wrapfig}
\makeatletter
\let\wrapped@colsep\relax
\newcommand*\wrapped@setcolsep[1]{%
  \ifx\wrapped@colsep\relax \xdef\wrapped@colsep{\the\columnsep}\fi
  \global\columnsep=\dimexpr #1\relax}
\newcommand*\wrapped@restorecolsep{%
  \ifx\wrapped@colsep\relax\else
    \global\columnsep=\wrapped@colsep\relax
    \global\let\wrapped@colsep\relax
  \fi}
\g@addto@macro\WF@startwrapping{\wrapped@restorecolsep}
% \WFclear as in wrapfig, but a figure that waits is centered. LaTeX runs
% \WFclear at each \clearpage, the end of the document included.
\renewcommand\WFclear{\par
  \ifvoid\WF@box\else
    \vskip\intextsep
    \hbox to\linewidth{\hss\box\WF@box\hss}%
    \vskip\intextsep
    \let\everypar\WF@@everypar \everypar\expandafter{\the\WF@everypar}%
  \fi
  \global\c@WF@wrappedlines\z@ \WF@finale
  \wrapped@restorecolsep}
\newcommand*\wrapped@endwrap{%
  \ifvoid\WF@box\else \WFclear \fi
  \ifx\parshape\WF@fudgeparshape
    \par
    \@tempcnta\c@WF@wrappedlines \advance\@tempcnta\m@ne
    \vspace{\@tempcnta\baselineskip}%
    \WFclear
  \fi}
% wrapfig runs \WF@startfloating after the environment's groups end. Before
% it, \wrapped@afterenv ends an earlier wrap there and records that group
% level for \WF@floathand below.
\def\wrapped@level{0}
\newcommand*\wrapped@afterenv{%
  \xdef\wrapped@level{\the\currentgrouplevel}%
  \ifx\parshape\WF@fudgeparshape \WF@finale \fi}
\let\wrapped@startfloating\WF@startfloating
\def\WF@startfloating{\wrapped@afterenv\wrapped@startfloating}
% A paragraph inside a group, such as a heading or the end of a callout,
% can start where a moved figure waits. A figure that wrapfig placed there
% would lose its wrap when the group ends, and the text after it would run
% over the figure. So wrapfig places a moved figure only at a paragraph on
% the group level where its environment ended.
\let\wrapped@floathand\WF@floathand
\def\WF@floathand{%
  \ifnum\currentgrouplevel=\wrapped@level\relax
    \expandafter\wrapped@floathand
  \fi}
\makeatother
\ExplSyntaxOn
\tl_new:N \l__wrapped_side_tl
\keys_define:nn { wrapped }
  {
    side .choice:,
    side / left .code:n = \tl_set:Nn \l__wrapped_side_tl { L },
    side / right .code:n = \tl_set:Nn \l__wrapped_side_tl { R },
    side .initial:n = right,
    width .tl_set:N = \l__wrapped_width_tl,
    width .initial:n = 0.4\linewidth,
    gap .tl_set:N = \l__wrapped_gap_tl,
    gap .initial:n = 1em,
    lines .tl_set:N = \l__wrapped_lines_tl,
    lines .initial:n = ,
  }
\NewDocumentEnvironment { wrappedfigure } { O{} }
  {
    \keys_set:nn { wrapped } { #1 }
    \use:c { wrapped@endwrap }
    % A caption in a narrow box reads better set plain and ragged-right.
    \cs_if_exist:NT \captionsetup
      { \captionsetup { format = plain, justification = raggedright } }
    \use:e
      {
        \exp_not:N \wrapfigure
        \tl_if_empty:NF \l__wrapped_lines_tl
          { [ \exp_not:V \l__wrapped_lines_tl ] }
        { \exp_not:V \l__wrapped_side_tl }
        { \exp_not:V \l__wrapped_width_tl }
      }
  }
  {
    \endwrapfigure
    \use:c { wrapped@setcolsep } { \l__wrapped_gap_tl }
  }
\ExplSyntaxOff
